\clearpage
\fancyhead[L]{\small VESTA.}
\hypertarget{native-295-lower}{}
\centerline{VESTA.}

\begin{center}
\rule{0.5\linewidth}{0.4pt}\\[2pt]
Monatliche Correspondenz. Vol. XVI. Pp. 285–291. September 1807.\\[2pt]
\rule{0.5\linewidth}{0.4pt}
\end{center}

\smallskip
\noindent---\ ---\ ---

\smallskip
\par\indent “I ended my observations on 8 July, Dr. \textsc{Olbers} writes to us, because the planet’s low altitude, moonlight, and twilight made them too difficult. The last six observations were made in the presence of our excellent friend \textsc{Gauss}, whose visit gave me pleasure. This excellent astronomer will soon give us the III. elements of \emph{Vesta}. The second elements, meanwhile, will need only slight corrections, since they still represent the latest observations so accurately, giving the right ascension only about $1'$ too large and the declination less than $30''$ too small. From the II. elements I have found that on 25 March 1796 \emph{Vesta} was in the zone observed at that time according to the \emph{Hist. Cél.} But a small and quite possible
\SourceLeaf{296}error in the semimajor axis, and therefore in \emph{Vesta}’s mean motion, could still change its declination for that day considerably. I do not know whether a star observed in that zone is now missing, because the part in which \emph{Vesta} could have been is now hidden in the Sun’s glare.”

\par\indent We received these \emph{third elements} a few days ago and give them below. In addition to all the observations already published in this journal, Dr. \textsc{Gauss} used a series of Paris observations from 13 April to 22 May that Mr. \textsc{Bouvard} had sent him.

\smallskip

\begin{center}\begin{tabular}{lr}
Epoch 31 March 1807, noon in Bremen  & $192^\circ\;23'\;30''1$ \\
Daily tropical motion .\ .\ . & $981''8459$ \\
\shortstack[l]{Aphelion, held fixed in sidereal\\coordinates for the epoch} .\ .\ .\ .\ .\ .\ .\ .\ .\ . & $69^\circ\;50'\;31''9$ \\
Ascending node .\ .\ .\ .\ . & $103^\circ\;18'\;28''$ \\
Inclination of the orbit .\ .\ .\ .\ .\ .\ . & $7\;8\;10.7$ \\
Eccentricity .\ .\ .\ .\ .\ .\ .\ .\ .\ . & $0.0855050$ \\
Logarithm of the mean distance .\ . & $0.3720160$ \\
\end{tabular}\end{center}
\par\indent The difference between these elements and the second is slight, yet the resulting geocentric position of \emph{Vesta} at the next opposition, in September 1808, will be $22'$ greater. Dr. \textsc{Gauss} was also kind enough to send us that series of Paris observations, along with the differences between the observed positions and those calculated from the II. elements.

\smallskip
\clearpage\begingroup\fontsize{9}{11}\selectfont\setlength{\tabcolsep}{2.5pt}\renewcommand{\arraystretch}{1.05}
\begin{center}\begin{tabular}{lrrrrr}
1807 & \shortstack{Mean time\\in Paris} & R.A. \VestaRATwoNinetySix & \shortstack{Diff. from\\II. elements} & Decl. \VestaDecTwoNinetySix & \shortstack{Diff. from\\II. elements} \\ \hline
April 13 & $22^{\text{u}}\;39^{\text{m}}\;6\mathord{\raisebox{.65ex}{\ObservedTimeStroke}}2$ & $180^\circ 57' 7''8$ & $+3''4$ & $12^\circ 50' 35''0$ & $+5''2$ \\
14 & $22\;34\;28.3$ & $180\;46\;27.9$ & $+1.3$ & $12\;52\;39.2$ & $+6.0$ \\
18 & $22\;16\;7.0$ & $180\;\;7\;\;6.9$ & $-3.1$ & $12\;58\;\;9.3$ & $+12.7$ \\
19 & $22\;11\;34.9$ & $179\;58\;\;4.1$ & $+2.2$ & $12\;59\;\;5.5$ & $+0.4$ \\
21 & $22\;2\;39.0$ & $179\;41\;20.6$ & $-1.4$ & $12\;58\;40.7$ & $+65.0$ \\
24 & $21\;49\;20\;2$ & $179\;19\;\;5.1$ & $+1.7$ & $12\;58\;44.2$ & $+1.7$ \\
25 & $21\;44\;58.3$ & $179\;12\;26.6$ & $+4.7$ & $12\;57\;47.3$ & $+7.4$ \\
26 & $21\;40\;38.5$ & $179\;\;6\;33.8$ & $-13.1$ & $12\;56\;40.6$ & $+7.6$ \\
27 & $21\;36\;18.4$ & $179\;\;0\;35.4$ & $+0.1$ & $12\;55\;10.8$ & $+15.9$ \\
28 & $21\;32\;1.2$ & $178\;55\;15.2$ & $-0.$ & $12\;53\;47.3$ & $+2.8$ \\
29 & $21\;27\;45.5$ & $178\;50\;16.5$ & $+4.1$ & $12\;51\;54.4$ & $+4.3$ \\
30 & $21\;23\;31.7$ & $178\;45\;52.5$ & $-0.8$ & $12\;49\;49.4$ & $+3.2$ \\
May 1 & $21\;19\;20.0$ & $178\;41\;47.7$ & $+1.6$ & $12\;47\;26.9$ & $+5.0$ \\
7 & $20\;54\;44.7$ & $178\;26\;43.7$ & $+4.3$ & $12\;28\;39.6$ & $-7.5$ \\
12 & $20\;35\;3.3$ & $178\;26\;18.3$ & $+8.1$ & $12\;\;6\;35.8$ & $+2.8$ \\
15 & $20\;23\;35.2$ & $178\;31\;18.5$ & $+7.9$ & $11\;50\;39.4$ & $+25.5$ \\
16 & $20\;19\;50.7$ & $178\;33\;47.4$ & $+9.8$ & $11\;45\;30.8$ & $+0.2$ \\
17 & $20\;16\;5.0$ & $178\;36\;42.3$ & $+10.9$ & $11\;39\;46.8$ & $-0.6$ \\
18 & $20\;13\;22.5$ & $178\;39\;57.3$ & $+16.6$ & $11\;33\;50.4$ & $+0.3$ \\
20 & $20\;5\;1.9$ & $178\;47\;55.5$ & $+13.0$ & $11\;21\;31.8$ & $-3.5$ \\
21 & $20\;1\;24.6$ & $178\;52\;29.4$ & $+12.4$ & $11\;15\;\;5.5$ & $-3.6$ \\
22 & $19\;57\;48.9$ & $178\;57\;25.8$ & $+12.9$ & $11\;\;8\;27.6$ & $-2.1$ \\
\end{tabular}\end{center}\endgroup


\par\indent The times here are reckoned from midnight, and the declination on 21 April is evidently one arcminute too small. But who could fail to marvel at this excellent agreement of the II. elements with the sky?

\par\indent “I enclose the ephemeris for the path of \emph{Vesta} next year, Dr. \textsc{Gauss} goes on to write to us. Much of \emph{Vesta}’s path on this occasion falls on \textsc{Harding}’s chart for the motion of \emph{Ceres} in 1804, which accompanies the M. Corr. Professor \textsc{Harding} will soon give the sheet of his atlas that covers the rest to the copper engraver. This useful undertaking cannot be recommended and supported enough. I have already seen the first, very neatly engraved sheet in \emph{Bremen}.”

\SourceLeaf{297}\par\indent Unfortunately, according to the accompanying ephemeris, \emph{Vesta} will only be visible again in May or June of the following year. [Path of Vesta from 2 May 1808 to 18 March 1809 at four-day intervals according to the III. elements. R.A. and declination to the arcminute, logarithm of the distance to four decimal places.]

\par\indent Now that the elements of \emph{Vesta}’s orbit are already fairly close, it will be interesting to try to find it in older star catalogues. Our effort to find an observation of \emph{Vesta} among \textsc{Mayer}’s observations has so far been unsuccessful.

\begin{flushright}
[\textsc{von Zach}.]
\end{flushright}

